\documentclass[11pt,a4paper,twoside]{article}

% =========================
% Packages
% =========================
\usepackage[margin=1in]{geometry}
\usepackage[T1]{fontenc}
\usepackage[utf8]{inputenc}
\usepackage{lmodern}
\usepackage{microtype}
\usepackage{setspace}
\usepackage{amsmath,amssymb,amsfonts}
\usepackage{graphicx}
\usepackage{float}
\usepackage{booktabs}
\usepackage{tabularx}
\usepackage{hyperref}
\usepackage{caption}
\usepackage{fancyhdr}

% =========================
% Formatting
% =========================
\setstretch{1.1}
\captionsetup{font=small,labelfont=bf}
\hypersetup{
    colorlinks=true,
    linkcolor=black,
    citecolor=black,
    urlcolor=blue
}
\setlength{\parindent}{1.2em}
\setlength{\parskip}{0.4em}
\raggedbottom

% =========================
% Header / Footer
% =========================
\pagestyle{fancy}
\fancyhf{}
\fancyhead[RO]{\small Predicting Product Review Helpfulness\hspace{1em}\thepage}
\fancyhead[LE]{\small\thepage\hspace{1em}Group 10 et al.}
\fancyfoot{}
\renewcommand{\headrulewidth}{0pt}
\renewcommand{\footrulewidth}{0pt}

\fancypagestyle{plain}{
    \fancyhf{}
    \fancyhead[RO]{\small Predicting Product Review Helpfulness\hspace{1em}\thepage}
    \fancyhead[LE]{\small\thepage\hspace{1em}Group 10 et al.}
    \fancyfoot{}
    \renewcommand{\headrulewidth}{0pt}
    \renewcommand{\footrulewidth}{0pt}
}

% =========================
% Title & Authors
% =========================
\title{\textbf{Predicting Product Review Helpfulness using Linguistic Analysis and Explainable AI (XAI)}}
\author{%
\begin{tabular}{c}
Group 10 - DAP391m\\[0.45em]
FPT University\\[0.25em]
\end{tabular}
}
\date{}

\begin{document}
\maketitle
\thispagestyle{empty}

% =========================
% Abstract
% =========================
\begin{abstract}
The proliferation of online reviews has made it difficult for consumers to identify truly helpful feedback. This paper presents a machine learning framework to predict review helpfulness using the Amazon Reviews 2023 (All Beauty) dataset. We extract structural, linguistic (POS, NER, Readability), and sentiment features, combined with TF-IDF n-grams, to train classifiers including Logistic Regression, Random Forest, XGBoost, LinearSVC, and ComplementNB. To handle class imbalance, we applied `scale_pos_weight` and threshold tuning, increasing Recall significantly. Furthermore, we integrated an Explainable AI (XAI) layer using Google Gemini to provide transparent decision support. Our results demonstrate that combining linguistic complexity with robust algorithms and generative AI effectively distinguishes helpful reviews and provides actionable insights.
\end{abstract}

\noindent\textbf{Keywords:} Review Helpfulness, Natural Language Processing (NLP), XGBoost, Explainable AI (XAI), Class Imbalance.

% =========================
% 1. Introduction
% =========================
\section{Introduction}
Online reviews are a cornerstone of modern e-commerce, influencing consumer trust and purchasing decisions. However, the sheer volume of reviews, many of which are brief or uninformative, creates an ``information overload'' problem \cite{amazon2023}. Predicting the ``helpfulness'' of a review—the likelihood that a community will find the review useful—is therefore a critical task for e-commerce platforms.

Existing helpfulness prediction systems often provide only numerical scores without explaining \textit{why} a review is helpful. This paper addresses that gap by presenting an end-to-end framework encompassing data processing, optimized anomaly-aware modeling (due to severe class imbalance), and GenAI-powered explanations.

% =========================
% 2. Related Work
% =========================
\section{Related Work}
Prior work in review helpfulness prediction typically focuses either purely on sentiment analysis or on basic metadata (e.g., star rating, length). Recent studies have shown that linguistic features, such as readability and Part-Of-Speech (POS) distributions, carry stronger signals for helpfulness \cite{flesch}. However, most implementations stop at providing a classification label. Our work extends these baselines by combining advanced feature engineering (SpaCy, TextStat) with Generative AI (Google Gemini) to transform raw predictions into natural language explanations.

% =========================
% 3. Proposed Methodology
% =========================
\section{Proposed Methodology}

\subsection{Data Collection and Preprocessing}
We utilized the Amazon Reviews 2023 dataset (All Beauty category). The raw data, containing ~50,000 records in JSONL format, was cleaned using regular expressions to normalize text and remove special characters. The helpfulness label was formulated as a binary threshold ($votes \geq 1$) addressing a severe positive class rarity (~1.7\%).

\subsection{Feature Engineering}
Our feature space merges sparse textual matrices with dense metadata:
\begin{itemize}
    \item \textbf{Textual}: TF-IDF vectorization with bigram support (ngram\_range=1,2, max\_features=2000).
    \item \textbf{Sentiment}: VADER Compound Polarity scores.
    \item \textbf{Linguistic (Advanced)}: Flesch Reading Ease, Named Entity count, Noun/Adjective ratios (via SpaCy), and Comparative keyword counts.
\end{itemize}

\subsection{Modeling and Generative XAI}
We evaluated Logistic Regression, Random Forest, XGBoost, LinearSVC, and ComplementNB. For XGBoost, LinearSVC, and Naive Bayes, hyperparameters and thresholds were tuned via cross-validation, targeting the F1-Score metric instead of Accuracy, and incorporating class-weighting bounds (like `scale_pos_weight`) to penalize minority-class errors.
To interpret the predictions, we introduced an XAI module that passes the review, predicted label, and probability score to the Google Gemini API, which generates a brief, structured justification for the classification.

% =========================
% 4. Experimental Setup
% =========================
\section{Experimental Setup}
The experiments and the web dashboard were developed using Python 3.9+. The machine learning pipeline utilizes `scikit-learn` and `xgboost`. Linguistic processing relies on `spacy` (`en_core_web_sm`) and `textstat`. The interactive dashboard is built with `streamlit` and `plotly`. The XAI layer uses `google-genai` (Gemini 2.5 Flash).

% =========================
% 5. Results and Discussion
% =========================
\section{Results and Discussion}
The models were evaluated based on their ability to retrieve the minority class (Helpful reviews). Accuracy is misleading in this context; hence we optimized for F1-Score and Recall.

\begin{table}[ht]
\centering
\caption{Model Performance Snapshot (Hold-out Test Set)}
\label{tab:performance}
\begin{tabular}{l c c c c}
\toprule
\textbf{Model} & \textbf{Accuracy} & \textbf{F1-Score} & \textbf{Precision} & \textbf{Recall} \\
\midrule
Logistic Regression & 73.41\% & 25.94\% & 59.16\% & 16.61\% \\
Random Forest & 72.56\% & 7.87\% & 66.86\% & 4.18\% \\
XGBoost (Optimized) & 62.53\% & 49.99\% & 39.94\% & 66.81\% \\
LinearSVC & 56.01\% & 48.09\% & 35.93\% & 72.67\% \\
ComplementNB & 58.86\% & 48.59\% & 37.40\% & 69.35\% \\
\bottomrule
\end{tabular}
\end{table}

As shown in Table~\ref{tab:performance}, LinearSVC achieved the highest Recall at nearly 73\%, followed by ComplementNB and the optimized XGBoost, all representing massive improvements over standard Random Forest, effectively identifying helpful signals despite the data imbalance. Feature Importance analysis confirmed that VADER sentiment and Reading Ease played significant roles alongside standard review length.

% =========================
% 6. Conclusion
% =========================
\section{Conclusion}
This study demonstrates that review helpfulness can be effectively predicted using a hybrid of NLP features and gradient boosting designed for imbalanced data. More importantly, integrating LLMs (Google Gemini) elevates the predictive model from a ``black box'' to a transparent advisory tool. Future work will involve exploring in-house fine-tuned transformer architectures (e.g., RoBERTa) and scaling across varied review categories.

% =========================
% References
% =========================
\begin{thebibliography}{9}
\bibitem{amazon2023}
Hou, Y., et al. (2024). \textit{Amazon Reviews 2023}. Data repository.

\bibitem{flesch}
Kincaid, J. P., et al. (1975). \textit{Derivation of New Readability Formulas}. Technical Report.
\end{thebibliography}

\end{document}
